%% 自動生成（phys/gen_phys.py）
\phtocgrid{\phgi{保存則と相似の地図}{4}\phgi{この本の使い方}{6}}
{\color{bkLine}\hrule height 0.35pt}\vspace{2.4mm}
\phtocchap{電-0}{発電所の電気は、どうやって家まで届く？}{送電線は、なぜとても高い電圧で電気を送る？}{}{3}{3}{8}{J}
\phtocpart{第1部}{引きあう：静電場と電流}{12}{1}
\phtocchap{電-1}{電気の力と場：クーロンの法則と電場}{下じきでこすった髪の毛は、なぜ立ち上がる？}{4,8}{4}{4}{14}{C}
\phtocchap{電-2}{電気力線を数える：ガウスの法則と電位}{電気の力は、なぜはなれるほど弱くなる？}{8}{5}{5}{30}{C}
\phtocchap{電-3}{電気をためる：導体・誘電体とコンデンサ}{雷が落ちても、金属の車の中の人が無事なのはなぜ？}{8}{4}{5}{48}{C}
\phtocchap{電-4}{電気の流れと熱：電流と抵抗}{ドライヤーは、電気でどうやって熱をつくる？}{8}{4}{4}{66}{C}
\phtocchap{電-5}{回路を行列で解く：直流回路の解析}{鳥は、なぜ電線にとまっても感電しない？}{3,7}{4}{4}{82}{C}
\phtocpart{第2部}{うず巻く：磁場と電磁誘導}{98}{2}
\phtocchap{電-6}{電流が方位磁針を動かす：電流がつくる磁場}{電線に電気を流すと、なぜ方位磁針がふれる？}{4,8}{5}{5}{100}{C}
\phtocchap{電-7}{鉄と磁石：磁性体と磁気回路}{磁石にくっついたくぎは、なぜ自分も磁石になる？}{5}{4}{5}{118}{C}
\phtocchap{電-8}{電気が生まれる：電磁誘導とインダクタンス}{IH調理器は、火がないのになぜなべが熱くなる？}{8}{5}{5}{136}{C}
\phtocbreak
\phtocpartcont{第2部}{うず巻く：磁場と電磁誘導}{2}
\phtocchap{電-9}{スイッチを入れた瞬間：過渡現象}{充電器をコンセントからぬいても、ランプがしばらく光っているのはなぜ？}{6,9}{5}{5}{154}{C}
\phtocpart{第3部}{変える：交流と電気機器}{172}{3}
\phtocchap{電-10}{向きが変わる電気：交流回路とフェーザ}{ラジオは、たくさんの電波からどうやって1つの局を選ぶ？}{5,6}{4}{4}{174}{C}
\phtocchap{電-11}{電力を運ぶ3本の線：交流電力と三相交流}{送電線の鉄塔では、なぜ電線が3本ずつ組になっている？}{4,6}{4}{5}{190}{C}
\phtocchap{電-12}{モータと変圧器のしくみ：電気機器}{電気自動車のモータは、どうやって回る？}{6}{4}{5}{208}{C}
\phtocchap{電-13}{長い電線を読む：二端子対回路と分布定数回路}{テレビのアンテナ線に、なぜふつうの電線を使わない？}{6,7}{5}{5}{226}{C}
\phtocpart{第4部}{届く：電磁波}{244}{4}
\phtocchap{電-14}{4つの式にまとめる：マクスウェル方程式}{電気と磁気は、別々のものなの？}{8}{5}{5}{246}{C}
\phtocchap{電-15}{光も電波も同じ波：電磁波}{電子レンジの窓の網は、何のためにある？}{8}{5}{5}{264}{C}
\phtocchap{電-16}{電気を届ける・電気を変える}{停電しないように、電気の量をどうやって合わせている？}{}{4}{4}{282}{R}
\par\vspace{1.5mm}
\phtocgrid{\phgi{こたえ}{290}\phgi{記号カード}{298}\phgi{物性値・定数・単位}{301}\phgi{\mbox{なぜ？}の道すじ}{308}\phgi{できるようになったことチェック表}{310}\phgi{つぎの巻へ}{312}\phgi{さくいん}{314}}
